% ---------------------------------------------------------------------
%  Chapitre 2 — Les données et leur structure de dépendance
%  (version aperçu, intégralement rédigée)
% ---------------------------------------------------------------------
\chapter{Les données et leur structure de dépendance}
\label{chap:donnees}

\section{La base des trades}
\label{sec:donnees-trades}

Tout au long de ce rapport, on appliquera les méthodes à notre base de
données. Celle-ci a été obtenue en appliquant différentes stratégies à
des données antérieures. Cette méthode se nomme le \emph{backtesting}~:
elle consiste à rejouer le passé pour vérifier si une stratégie a repéré
un signal de prédiction du marché.

Ici, nous appliquons dix stratégies techniques classiques ainsi qu'une
stratégie témoin, appelée oracle. La stratégie oracle utilise les
données futures pour effectuer des trades uniquement gagnants, afin de
vérifier que notre protocole laisse passer une stratégie dont le signal
fonctionne réellement. Les dix autres stratégies, elles, effectuent
leurs trades à une date fixée en utilisant uniquement les données
précédant cette date.

\begin{table}[H]
\centering
\caption{Les onze stratégies du banc d'essai et l'indicateur sur lequel
chacune repose.}
\label{tab:strategies}
\renewcommand{\arraystretch}{1.25}
\small
\begin{tabular}{@{}ll@{}}
\toprule
Stratégie & Fondée sur \\
\midrule
\texttt{rsi\_classic}  & indicateur RSI (seuils 30/70) \\
\texttt{rsi\_strict}   & indicateur RSI (seuils 20/80) \\
\texttt{rsi\_trend}    & indicateur RSI avec filtre de tendance (MM200) \\
\texttt{ma\_crossover} & croisement de deux moyennes mobiles \\
\texttt{db\_bottom}    & figure en double creux \\
\texttt{dt\_top}       & figure en double sommet \\
\texttt{hs\_inverse}   & figure épaule-tête-épaule inversée \\
\texttt{hs\_classic}   & figure épaule-tête-épaule \\
\texttt{sr\_breakout}  & cassure d'une résistance \\
\texttt{sr\_breakdown} & cassure d'un support \\
\midrule
\texttt{oracle}        & témoin~: connaît le cours futur (triche) \\
\bottomrule
\end{tabular}
\end{table}

J'ai alors choisi d'appliquer ces stratégies aux cinq cent une grandes
capitalisations américaines présentes dans le S\&P~500. Cet indice
constitue un ensemble sur lequel on peut investir avec peu de risque~:
même si le cours de chacune de ces actions peut chuter individuellement,
l'indice lui-même reste relativement stable et suit une tendance
haussière.

Pour qu'une stratégie soit validée, il faut qu'elle obtienne de
meilleurs résultats qu'une stratégie jouant au hasard, et que cela ne
soit pas dû à notre échantillonnage. Or la tendance est ici haussière,
ce qui nous ramène à la première difficulté de l'introduction~: comment
évaluer une stratégie lorsque le marché monte de lui-même~?

Pour répondre à cette question, on simule le hasard. À chaque trade
réalisé, on associe deux cents trades fictifs de même exposition, sur le
même titre, avec la même durée et le même sens, mais dont la date
d'entrée est tirée au sort. La moyenne de leurs rendements définit le
gain du hasard pour ce trade. On définit alors, pour chaque trade, son
\edgevar{} comme le rendement effectivement obtenu par la stratégie,
éventuellement négatif, diminué de ce gain du hasard~:
\[
  \edgevar = \text{gain de la stratégie} - \text{gain du hasard}.
\]
Un \edgevar{} positif signifie que la stratégie a fait mieux qu'une
entrée au hasard de même exposition~: son \emph{timing} apporte une
information que la seule dérive du marché n'explique pas.

On obtient ainsi une base contenant, pour chaque stratégie, entre 978 et
8\,513 trades effectués entre janvier~2010 et juillet~2026, pour un
total de 40\,362 trades. Chaque ligne correspond à une position ouverte
puis fermée par une stratégie sur un titre. On y associe le titre
concerné, désigné par son \emph{ticker}, c'est-à-dire son code en
bourse, la stratégie employée, le secteur d'activité, les dates d'entrée
et de sortie, la durée de détention, le régime de marché et le contexte
à l'entrée, le rendement obtenu et l'avantage \edgevar{} sur un tirage
aléatoire équivalent.

\begin{table}[H]
\centering
\caption{Description des variables d'un trade, illustrées sur un trade réel
de la base (achat du titre \texttt{A} par la stratégie \texttt{rsi\_classic}).}
\label{tab:dico-trades}
\renewcommand{\arraystretch}{1.3}
\small
\begin{tabular}{@{}>{\raggedright\arraybackslash}p{2.4cm}>{\raggedright\arraybackslash}p{2.6cm}>{\raggedright\arraybackslash}p{5.6cm}l@{}}
\toprule
Variable & Nature & Signification & Exemple \\
\midrule
\texttt{ticker}       & qualitative (501 mod.) & le titre traité (une même action revient dans plusieurs trades) & \texttt{A} \\
\texttt{strategy}     & qualitative (11 mod.)  & la stratégie ayant produit le trade & \texttt{rsi\_classic} \\
\texttt{sector}       & qualitative (11 mod.)  & le secteur d'activité du titre & Santé \\
\texttt{entry\_date}  & date                   & date d'achat & 2011-08-11 \\
\texttt{exit\_date}   & date                   & date de revente & 2012-01-20 \\
\texttt{holding\_days}& quantitative           & durée de détention, en jours & $162$ \\
\texttt{regime\_entry}& qualitative            & régime du marché à l'entrée & baissier \\
\texttt{vol\_entry}   & quantitative           & volatilité du titre à l'entrée & $0{,}043$ \\
\texttt{rsi\_entry}   & quantitative           & niveau de l'indicateur RSI à l'entrée & $30{,}1$ \\
\texttt{dist\_ma200}  & quantitative           & écart à la moyenne mobile sur 200 jours & $-0{,}210$ \\
\texttt{return\_net}  & quantitative           & rendement net réalisé par la stratégie & $+0{,}187$ \\
\texttt{rand\_return} & quantitative           & rendement d'une entrée au hasard équivalente & $+0{,}078$ \\
\texttt{edge}         & quantitative           & avantage $=$ \texttt{return\_net} $-$ \texttt{rand\_return} & $+0{,}109$ \\
\texttt{win}          & binaire                & trade gagnant ($1$) ou perdant ($0$) & $1$ \\
\bottomrule
\end{tabular}
\end{table}

\section{Les rendements sectoriels}
\label{sec:donnees-secteurs}

Pour étudier la dépendance des données, on a besoin de données par
secteur. Le S\&P~500 compte onze secteurs d'activité. Sur chacun d'eux,
on établit une courbe du rendement journalier, obtenu en moyennant
chaque jour les rendements des titres du secteur. Ces données seront
utilisées d'une part au canal~2 du paragraphe~\ref{sec:dep-canal2}, pour
mettre en évidence le facteur marché, c'est-à-dire la part des
mouvements commune à tous les secteurs, et d'autre part dans la
conclusion du chapitre~\ref{chap:juger}, pour examiner si le marché est
prévisible à partir de son propre passé.

\section{La structure de dépendance, mesurée}
\label{sec:dep-mesuree}

Les tests que nous utiliserons au chapitre~\ref{chap:juger} ne
fonctionnent que si les données sont indépendantes. Ici, les données
sont dépendantes par le titre, par la date et par la durée. Cette
situation porte un nom en théorie des sondages~: notre base n'est pas un
tirage au hasard d'individus, mais un tirage par grappes, où chaque
grappe est un titre. Le réflexe du sondeur impose alors de décrire ce
plan d'échantillonnage avant d'estimer quoi que ce soit. C'est pourquoi
on mesure les dépendances maintenant, avant les tests.

\subsection{Canal 1 --- la dépendance de titre}
\label{sec:dep-canal1}

Deux trades d'un même titre sont liés par la gestion de la société ainsi
que par son image à l'extérieur. Ils subissent les mêmes chocs propres à
la société, ses résultats, ses décisions de direction, ses litiges,
indépendamment du moment où ils sont ouverts. On appellera cet effet la
dépendance de titre~: nos données constituent alors un tirage par
grappes, où chaque grappe est un titre.

On définit alors l'estimateur suivant~:
\[
  \hat\rho \;=\; \frac{\mathrm{MSB}-\mathrm{MSW}}
                      {\mathrm{MSB} + (n_0 - 1)\,\mathrm{MSW}},
\]
où MSB (\emph{Mean Square Between}) est la variation entre les titres,
donc inter-grappes, c'est-à-dire la moyenne des écarts d'\edgevar{} d'un
titre à l'autre, et MSW (\emph{Mean Square Within}) la variation à
l'intérieur d'un titre, donc intra-grappe, c'est-à-dire la moyenne des
écarts d'\edgevar{} entre les trades d'un même titre. Deux cas se
présentent. Si les trades d'un même titre ne se ressemblent pas plus que
ceux de titres différents, la variation est la même à l'intérieur des
titres et entre eux, et l'on obtient MSB~$\approx$~MSW. S'ils se
ressemblent, c'est la dépendance de titre~: la variation à l'intérieur
d'un titre s'effondre tandis que la variation entre titres subsiste, et
l'on obtient MSB~$>$~MSW.

On pourrait se contenter du quotient de MSB par MSW, qui vaudrait un en
cas d'indépendance. Mais un écart à un, qu'il soit de un dixième, de
trois ou de cent, ne se rattache à aucune échelle, et l'on ne saurait
dire s'il traduit une dépendance forte ou faible. Pour obtenir une
réponse du type « treize pour cent du résultat est dû à une dépendance
de grappe », on préfère la formule ci-dessus, qui ramène la mesure à une
valeur comprise entre zéro et un.

Le terme $n_0$ correspond à la taille moyenne des grappes, c'est-à-dire
au nombre de trades par titre, corrigée de l'inégalité des tailles. Il
pondère la variation intra-titre au dénominateur~: plus les grappes
contiennent de trades, plus $n_0$ est grand, et plus il faut de
contraste entre les titres pour conclure à une dépendance.

Le tableau~\ref{tab:dep-mesures} donne un $\hat\rho$ quasi nul pour la
plupart des stratégies, mais atteignant $0{,}134$ pour
\texttt{rsi\_strict}. La dépendance par le titre est donc négligeable
pour toutes les stratégies, excepté \texttt{rsi\_strict}. Un $\hat\rho$
faible ne garantit toutefois pas l'indépendance des trades~: il n'écarte
que ce premier canal, les deux suivants restant à examiner.

\subsection{Canal 2 --- la dépendance de marché}
\label{sec:dep-canal2}

La deuxième dépendance que l'on va mesurer est la dépendance du moment.
Des trades ouverts au même moment peuvent être liés parce qu'ils
subissent la même journée de bourse~: à une même période, différentes
actions changent de valeur conjointement, et les trades effectués sont
alors liés par le marché.

Cette dépendance ne s'observe pas directement sur les trades, mais sur
le marché lui-même. Si les secteurs bougent ensemble, alors des trades
ouverts au même moment, quel que soit leur secteur, subissent la même
poussée. Pour la mesurer, on utilise donc les données introduites au
paragraphe~\ref{sec:donnees-secteurs}, c'est-à-dire les rendements
quotidiens par secteur. On calcule les corrélations des secteurs deux à
deux, puis on en fait la moyenne, et l'on obtient une corrélation de
$0{,}70$. Ce chiffre élevé signifie que les secteurs sont largement
corrélés. Lorsqu'un secteur bouge, le marché a tendance à bouger dans le
même sens, et inversement, lorsque le marché suit une tendance, la
plupart des secteurs la suivent aussi.

On utilise ensuite la méthode de l'ACP sur les onze secteurs, de façon à
mettre en évidence les plus grandes corrélations
(figure~\ref{fig:acp-marche}). Cette méthode calcule les axes qui
captent la plus grande variance. En l'appliquant ici, on observe un
premier axe qui capte 73~\% de la variance totale, le long duquel tous
les secteurs progressent ensemble~: c'est le mouvement global du marché.
Vient ensuite un deuxième axe, nettement moins important, qui sépare les
secteurs cycliques, réagissant fortement au marché, des secteurs
défensifs, qui ne bougent guère quelle que soit la conjoncture. On peut
donner en exemple de secteur cyclique la technologie, qui dépend du
pouvoir d'achat, on achète des biens technologiques quand on est aisé,
moins en temps de crise, et comme secteur défensif les services aux
collectivités (électricité, eau, gaz), dont la consommation reste
constante, que le marché soit haussier ou en crise.

\begin{figure}[H]
  \centering
  \includegraphics[width=0.78\linewidth]{etape3_acp.png}
  \caption{ACP des onze rendements sectoriels~: variance expliquée par
  composante (gauche) et position des secteurs sur les deux premiers
  axes (droite).}
  \label{fig:acp-marche}
\end{figure}

Sur l'histogramme de gauche, la première barre écrase toutes les
autres~: elle capte à elle seule 73~\% de la variance, contre moins de
8~\% pour la deuxième. Sur le nuage de droite, tous les secteurs se
situent du même côté de l'axe horizontal, ce qui confirme qu'ils
progressent ensemble, tandis que le second axe les répartit entre
défensifs et cycliques.

Cette cause commune n'est donc pas une hypothèse mais une mesure~: 73~\%
des mouvements sont portés par une force unique. Deux trades ouverts au
même moment subissent largement le même sort, ce qui les rend dépendants
du marché.

\subsection{Canal 3 --- la dépendance temporelle}
\label{sec:dep-canal3}

Dans la suite du rapport, nous aurons besoin d'agréger les trades par
mois. Or les trades durent~: en moyenne 132~jours, avec une médiane de
65~jours, et plus de 75~\% d'entre eux dépassent un mois. Un même trade,
encore ouvert, est donc compté dans plusieurs mois consécutifs. Deux
mois voisins partagent ainsi des trades communs et se ressemblent. Même
agrégées par mois, les observations ne sont donc pas indépendantes, et
c'est cette ressemblance d'un mois à l'autre qu'il faut mesurer.

Pour cela, on utilise l'autocorrélation d'ordre un, notée
$\hat\rho(1)$~: on décale la série mensuelle d'un mois et on corrèle la
série avec elle-même ainsi décalée. Elle mesure à quel point un mois
ressemble au précédent. On obtient les valeurs suivantes.

\begin{table}[H]
\centering
\caption{Les dépendances mesurées, stratégie par stratégie~:
corrélation intra-titre $\hat\rho$ (canal~1) et autocorrélation
mensuelle $\hat\rho(1)$ (canal~3).}
\label{tab:dep-mesures}
\small
\begin{tabular}{@{}lrr@{}}
\toprule
Stratégie & $\hat\rho$ intra-titre & $\hat\rho(1)$ mensuelle \\
\midrule
\texttt{db\_bottom}    & $0{,}000$ & $0{,}36$ \\
\texttt{dt\_top}       & $0{,}000$ & $0{,}39$ \\
\texttt{hs\_classic}   & $0{,}087$ & $0{,}15$ \\
\texttt{hs\_inverse}   & $0{,}001$ & $0{,}31$ \\
\texttt{ma\_crossover} & $0{,}003$ & $0{,}56$ \\
\texttt{rsi\_classic}  & $0{,}000$ & $0{,}23$ \\
\texttt{rsi\_strict}   & $0{,}134$ & $0{,}03$ \\
\texttt{rsi\_trend}    & $0{,}010$ & $0{,}20$ \\
\texttt{sr\_breakdown} & $0{,}051$ & $0{,}11$ \\
\texttt{sr\_breakout}  & $0{,}004$ & $0{,}05$ \\
\bottomrule
\end{tabular}
\end{table}

Le $\hat\rho(1)$ atteint $0{,}56$ pour \texttt{ma\_crossover}, sa valeur
la plus forte, tandis qu'il reste faible pour \texttt{rsi\_strict}
($0{,}03$). La dépendance temporelle est a priori forte pour les
stratégies à trades longs, faible pour celles à trades courts.
Attention~: ces valeurs élevées ne traduisent pas une prévisibilité du
marché, mais le chevauchement mécanique des trades qui durent.

Les trois canaux de dépendance sont désormais mesurés. Aucun n'est
systématiquement élevé. La dépendance de titre reste souvent proche de
zéro, et même l'autocorrélation temporelle s'efface pour les stratégies
à trades courts. Mais chacun se manifeste fortement pour certaines
stratégies, la dépendance de titre pour \texttt{rsi\_strict}, la
dépendance de marché pour toutes, la dépendance temporelle pour
\texttt{ma\_crossover}, et le chapitre~\ref{chap:juger} montrera qu'ils
suffisent à renverser des verdicts. On peut maintenant juger les
stratégies en sachant exactement ce qui menace les tests.
